\documentclass[12pt, a4paper]{article}

% ---------- Encoding & Language ----------
\usepackage[utf8]{inputenc}
\usepackage[T1]{fontenc}
\usepackage[ngerman]{babel}

% ---------- Fonts ----------
\usepackage{lmodern}

% ---------- Layout ----------
\usepackage[
  top=2.5cm,
  bottom=2.5cm,
  left=3cm,
  right=2.5cm
]{geometry}
\setlength{\parskip}{6pt}
\setlength{\parindent}{0pt}

% ---------- Header / Footer ----------
\usepackage{fancyhdr}
\pagestyle{fancy}
\fancyhf{}
\fancyhead[L]{\small\leftmark}
\fancyhead[R]{\small\thepage}
\renewcommand{\headrulewidth}{0.4pt}

% ---------- Sections styling ----------
\usepackage{titlesec}

% Section: groß, fett, mit Linie darunter
\titleformat{\section}
  {\large\bfseries}
  {\thesection}{1em}{}
  [\vspace{2pt}\hrule height 0.8pt\vspace{4pt}]
\titlespacing{\section}{0pt}{14pt}{6pt}

% Subsection: normal fett, deutlich sichtbar
\titleformat{\subsection}
  {\normalsize\bfseries}
  {\thesubsection}{1em}{}
  [\vspace{1pt}\hrule height 0.3pt\vspace{2pt}]
\titlespacing{\subsection}{0pt}{10pt}{4pt}

% Subsubsection: normalsize, fett
\titleformat{\subsubsection}
  {\normalsize\bfseries}
  {\thesubsubsection}{1em}{}
\titlespacing{\subsubsection}{0pt}{8pt}{2pt}

% ---------- Boxes ----------
\usepackage[most]{tcolorbox}

\newtcolorbox{definitionbox}[1][]{
  colback=white, colframe=black, coltext=black,
  fonttitle=\bfseries, title={Definition}, #1
}
\newtcolorbox{notebox}[1][]{
  colback=white, colframe=black, coltext=black,
  fonttitle=\bfseries, title={Hinweis}, #1
}
\newtcolorbox{examplebox}[1][]{
  colback=white, colframe=black, coltext=black,
  fonttitle=\bfseries, title={Beispiel}, #1
}

% ---------- TikZ ----------
\usepackage{tikz}
\usetikzlibrary{arrows.meta, positioning, shapes.geometric, fit, backgrounds, decorations.pathreplacing}

% ---------- Math ----------
\usepackage{amsmath, amssymb}

% ---------- Tabellen ----------
\usepackage{booktabs}
\usepackage{multirow}
\usepackage{array}

% ---------- Quelltext ----------
\usepackage{listings}
\lstset{
  basicstyle=\ttfamily\small,
  keywordstyle=\bfseries,
  columns=fullflexible,
  breaklines=true,
  frame=single,
  framesep=4pt,
  xleftmargin=4pt,
  showstringspaces=false
}

% ---------- Lists ----------
\usepackage{enumitem}
\setlist[itemize]{leftmargin=1.5em, itemsep=2pt}
\setlist[enumerate]{leftmargin=1.5em, itemsep=2pt}

% ---------- Figures ----------
\usepackage{float}

% ---------- Hyperlinks ----------
\usepackage[hidelinks]{hyperref}

% ---------- Misc ----------
\usepackage{microtype}

% =========================================
%                DOCUMENT
% =========================================
\begin{document}

\section{Aussagen und Mengen}

\subsection{Aussagenlogik}

Eine Aussage ist ein Satz, dem genau einer der Wahrheitswerte wahr oder
falsch zukommt. Die Junktoren werden durch Wahrheitstafeln definiert, nicht
durch Umgangssprache; das ist der einzige Weg, Streitfälle zu entscheiden.

\begin{center}
\begin{tabular}{cc|cccc}
\toprule
$p$ & $q$ & $p \wedge q$ & $p \vee q$ & $p \to q$ & $p \leftrightarrow q$\\
\midrule
f & f & f & f & \textbf{w} & w\\
f & w & f & w & \textbf{w} & f\\
w & f & f & w & f & f\\
w & w & w & w & w & w\\
\bottomrule
\end{tabular}
\end{center}

\begin{notebox}
Die beiden fett gesetzten Zeilen sind die, über die gestritten wird: aus einer
falschen Prämisse folgt alles. Die Alternative wäre, dass die Aussage
"`jedes Element der leeren Menge ist rot"' falsch ist, und damit würde kein
Beweis über die leere Menge mehr funktionieren.
\end{notebox}

\subsubsection{Tautologie, Erfüllbarkeit, Äquivalenz}

\begin{definitionbox}
Eine Formel heißt \emph{Tautologie}, wenn sie unter jeder Belegung wahr ist,
und \emph{erfüllbar}, wenn sie unter mindestens einer Belegung wahr ist. Zwei
Formeln sind \emph{äquivalent}, wenn ihre Wahrheitstafeln übereinstimmen.
\end{definitionbox}

Wichtige Äquivalenzen:
\begin{align*}
  \neg(p \wedge q) &\equiv \neg p \vee \neg q & \text{(de Morgan)}\\
  \neg(p \vee q) &\equiv \neg p \wedge \neg q & \text{(de Morgan)}\\
  p \to q &\equiv \neg q \to \neg p & \text{(Kontraposition)}\\
  p \to q &\not\equiv q \to p & \text{(die Umkehrung ist \emph{nicht} äquivalent)}
\end{align*}

\subsubsection{Zwei Schlussregeln, die sich ähneln}

\begin{center}
\begin{tabular}{lll}
\toprule
Regel & Form & gültig\\
\midrule
Modus ponens & aus $p$ und $p \to q$ folgt $q$ & ja\\
Bejahung des Konsequens & aus $q$ und $p \to q$ folgt $p$ & \textbf{nein}\\
\bottomrule
\end{tabular}
\end{center}

Die zweite Regel wird durch eine einzige Zeile der Wahrheitstafel widerlegt:
$p$ falsch, $q$ wahr. Der ganze Beweisbegriff der Vorlesung beruht darauf,
diese beiden Regeln mechanisch unterscheiden zu können.

\subsubsection{Quantoren}

\begin{align*}
  \neg \forall x \in M: P(x) &\equiv \exists x \in M: \neg P(x)\\
  \neg \exists x \in M: P(x) &\equiv \forall x \in M: \neg P(x)
\end{align*}

Über der leeren Menge ist jede Allaussage wahr und jede Existenzaussage
falsch. Das ist keine Sonderregel, sondern folgt aus der Definition, und es
ist der Grund, warum in Beweisen der Fall $M = \emptyset$ meist ohne Arbeit
erledigt ist.

\subsection{Mengen}

\begin{definitionbox}
Für Mengen $A, B$ über einem Universum $U$:
\[
  A \cup B, \quad A \cap B, \quad A \setminus B, \quad
  A \triangle B = (A \cup B) \setminus (A \cap B), \quad
  \overline{A} = U \setminus A.
\]
Die \emph{Potenzmenge} $\mathcal{P}(A)$ ist die Menge aller Teilmengen von
$A$; es gilt $|\mathcal{P}(A)| = 2^{|A|}$.
\end{definitionbox}

\subsubsection{Gesetze}

Die Gesetze von de Morgan gelten für Mengen wie für Aussagen:
\[
  \overline{A \cup B} = \overline{A} \cap \overline{B}, \qquad
  \overline{A \cap B} = \overline{A} \cup \overline{B}.
\]

\begin{notebox}
Im Begleitcode wurden diese Gesetze über einem Universum mit sechs Elementen
für alle $64 \cdot 64 = 4096$ Paare von Teilmengen geprüft. Das ist ein
Beweis für dieses Universum und ein Experiment für die allgemeine Aussage.
Die Unterscheidung ist wichtig: eine Aufzählung kann das Gesetz nicht
beweisen, aber sie würde ein falsches sofort widerlegen.
\end{notebox}

\subsubsection{Inklusion und Exklusion}

\[
  |A_1 \cup \dots \cup A_n| =
  \sum_{\emptyset \neq I \subseteq \{1,\dots,n\}} (-1)^{|I|+1}
  \Bigl| \bigcap_{i \in I} A_i \Bigr|
\]

Das Summieren der Größen zählt jedes Element so oft, wie es in Mengen liegt.
Das Abziehen der paarweisen Schnitte korrigiert das, entfernt die
Dreifachschnitte aber zu oft, und so weiter.

\subsection{Abzählbarkeit}

\begin{definitionbox}
Eine Menge heißt \emph{abzählbar}, wenn es eine surjektive Abbildung
$\mathbb{N} \to M$ gibt, also wenn ihre Elemente in einer (eventuell
unendlichen) Liste aufgeführt werden können.
\end{definitionbox}

\subsubsection{Was aufgezählt werden kann}

Die Cantorsche Paarungsfunktion
\[
  \pi(x, y) = \frac{(x+y)(x+y+1)}{2} + y
\]
ist eine Bijektion $\mathbb{N}^2 \to \mathbb{N}$. Damit sind $\mathbb{N}^2$
und, über die Darstellung als Bruch, auch $\mathbb{Q}$ abzählbar. Eine
Aufzählung muss jedes Element irgendwann erreichen, nicht der Größe nach.

\subsubsection{Das Diagonalverfahren}

\begin{definitionbox}
Sei $(s_i)_{i \in \mathbb{N}}$ eine beliebige Folge unendlicher
Ziffernfolgen. Die Folge $d$ mit $d_i = (s_i)_i + 1 \bmod 10$ unterscheidet
sich von jedem $s_i$ an der Stelle $i$. Also ist keine Aufzählung
vollständig, und die Menge der unendlichen Folgen ist überabzählbar.
\end{definitionbox}

Das Argument setzt nichts über die Herkunft der Liste voraus. Genau deshalb
ist es ein Beweis und kein Einwand gegen einen bestimmten Versuch.

\subsubsection{Der Satz von Cantor}

Für jede Menge $M$ und jede Abbildung $f: M \to \mathcal{P}(M)$ ist
\[
  D = \{ x \in M \mid x \notin f(x) \}
\]
kein Bild: wäre $D = f(a)$, so gälte $a \in D \iff a \notin D$. Also ist
$\mathcal{P}(M)$ echt größer als $M$, für endliche wie für unendliche Mengen.
Es ist dasselbe Diagonalargument in anderer Schreibweise.

\subsection{Beweisverfahren}

\begin{definitionbox}
\textbf{Direkter Beweis.} Aus den Prämissen wird die Konklusion durch eine
Kette gültiger Schlüsse gewonnen.

\textbf{Indirekter Beweis (Widerspruch).} Man nimmt $\neg B$ an und leitet
einen Widerspruch her; wegen $\neg B \to \text{falsch}$ folgt $B$.

\textbf{Kontraposition.} Statt $A \to B$ zeigt man $\neg B \to \neg A$;
beide Formeln sind äquivalent.

\textbf{Fallunterscheidung.} Gilt $A_1 \vee A_2$ und folgt $B$ sowohl aus
$A_1$ als auch aus $A_2$, so gilt $B$.
\end{definitionbox}

\begin{examplebox}
\textbf{$\sqrt{2}$ ist irrational.} Angenommen $\sqrt{2} = p/q$ vollständig
gekürzt. Dann $p^2 = 2q^2$, also ist $p$ gerade, $p = 2k$, also
$4k^2 = 2q^2$, also $q^2 = 2k^2$, also ist auch $q$ gerade. Das widerspricht
der Teilerfremdheit. Der Beweis ist indirekt und benutzt zusätzlich, dass
$p^2$ gerade nur für gerades $p$ möglich ist, was selbst eine
Kontraposition ist.
\end{examplebox}

\section{Relationen und Funktionen}

\subsection{Eigenschaften von Relationen}

\begin{definitionbox}
Eine Relation $R \subseteq M \times M$ heißt
\begin{itemize}
  \item \emph{reflexiv}, wenn $(a,a) \in R$ für alle $a$,
  \item \emph{symmetrisch}, wenn aus $(a,b) \in R$ stets $(b,a) \in R$ folgt,
  \item \emph{antisymmetrisch}, wenn aus $(a,b), (b,a) \in R$ stets $a = b$ folgt,
  \item \emph{transitiv}, wenn aus $(a,b), (b,c) \in R$ stets $(a,c) \in R$ folgt.
\end{itemize}
\end{definitionbox}

\begin{examplebox}
Auf einer dreielementigen Menge gibt es $2^9 = 512$ Relationen. Durch
vollständige Aufzählung im Begleitcode:

\begin{center}
\begin{tabular}{lr}
\toprule
Eigenschaft & Anzahl\\
\midrule
alle & 512\\
reflexiv & 64\\
symmetrisch & 64\\
antisymmetrisch & 216\\
transitiv & 171\\
Äquivalenzrelationen & 5\\
Halbordnungen & 19\\
\bottomrule
\end{tabular}
\end{center}
\end{examplebox}

\begin{notebox}
Symmetrie und Antisymmetrie schließen einander nicht aus: die Gleichheit hat
beide Eigenschaften. Die leere Relation ist symmetrisch, antisymmetrisch und
transitiv und scheitert nur an der Reflexivität, weshalb diese separat
gefordert wird.
\end{notebox}

\subsection{Hüllen}

\begin{definitionbox}
Die \emph{Hülle} einer Relation bezüglich einer Eigenschaft ist die kleinste
Relation, die $R$ enthält und die Eigenschaft besitzt. Sie existiert, weil
die Eigenschaft unter Durchschnitt erhalten bleibt.
\end{definitionbox}

Die Reihenfolge ist nicht beliebig. Für $R = \{(1,2),(2,3),(3,4)\}$:
\begin{itemize}
  \item erst symmetrisch, dann transitiv, dann reflexiv: eine
        Äquivalenzrelation;
  \item erst transitiv, dann symmetrisch: das Ergebnis ist nicht einmal
        transitiv.
\end{itemize}

Deshalb ist die Äquivalenzhülle in dieser Reihenfolge definiert, und nicht
als "`die drei Hüllen in beliebiger Reihenfolge"'.

\subsection{Äquivalenzrelationen und Partitionen}

\begin{definitionbox}
Eine Äquivalenzrelation auf $M$ zerlegt $M$ in Klassen; umgekehrt liefert
jede Partition von $M$ eine Äquivalenzrelation. Beide Konstruktionen sind
zueinander invers.
\end{definitionbox}

Die Anzahl der Partitionen einer $n$-elementigen Menge ist die Bellsche Zahl:

\begin{center}
\begin{tabular}{lrrrrrr}
\toprule
$n$ & 0 & 1 & 2 & 3 & 4 & 5\\
\midrule
$B_n$ & 1 & 1 & 2 & 5 & 15 & 52\\
\bottomrule
\end{tabular}
\end{center}

Die Kongruenz modulo $n$ ist das Beispiel, das in die Algebra weiterführt:
alle ihre Klassen sind gleich groß, was für Äquivalenzrelationen im
Allgemeinen falsch ist, und genau diese Gleichmäßigkeit macht die
Faktorgruppe möglich.

\subsection{Funktionen}

\begin{definitionbox}
$f: A \to B$ heißt \emph{injektiv}, wenn $f(a) = f(a')$ nur für $a = a'$
gilt; \emph{surjektiv}, wenn jedes $b \in B$ getroffen wird; \emph{bijektiv},
wenn beides gilt.
\end{definitionbox}

Für $|A| = m$, $|B| = n$:
\begin{align*}
  \text{alle Abbildungen} &: n^m\\
  \text{injektive} &: n (n-1) \cdots (n-m+1)\\
  \text{surjektive} &: \sum_{k=0}^{n} (-1)^k \binom{n}{k} (n-k)^m\\
  \text{bijektive} &: n! \text{ falls } m = n, \text{ sonst } 0
\end{align*}

Die Formel für die surjektiven Abbildungen ist wieder Inklusion und
Exklusion, diesmal über die nicht getroffenen Werte.

\begin{notebox}
Das Schubfachprinzip ist der Spezialfall $m > n$: eine Abbildung in eine
kleinere Menge kann nicht injektiv sein. Im Begleitcode gibt die Funktion das
kollidierende Paar zurück, nicht nur einen Wahrheitswert, weil die
interessanten Anwendungen die beiden Objekte benennen.
\end{notebox}

\subsection{Komposition und Umkehrung}

Für $f: A \to B$ und $g: B \to C$ ist $(g \circ f)(a) = g(f(a))$. Es gilt:
\begin{itemize}
  \item sind $f$ und $g$ injektiv, so auch $g \circ f$;
  \item sind $f$ und $g$ surjektiv, so auch $g \circ f$;
  \item ist $g \circ f$ injektiv, so ist $f$ injektiv (aber nicht
        notwendig $g$);
  \item ist $g \circ f$ surjektiv, so ist $g$ surjektiv (aber nicht
        notwendig $f$).
\end{itemize}

Die letzten beiden Aussagen sind die, die man sich merken muss: aus einer
Eigenschaft der Komposition folgt jeweils nur die Eigenschaft \emph{einer}
der beiden Abbildungen.

Eine Umkehrabbildung $f^{-1}$ existiert genau dann, wenn $f$ bijektiv ist.
Für endliche Mengen gleicher Größe genügt Injektivität oder Surjektivität,
für unendliche nicht: $n \mapsto n+1$ auf $\mathbb{N}$ ist injektiv und
nicht surjektiv.

\section{Induktion}

\subsection{Induktives Definieren}

\begin{definitionbox}
Eine induktive Definition nennt Basiselemente und Regeln, die aus vorhandenen
Elementen neue erzeugen. Definiert ist die \emph{kleinste} Menge, die die
Basiselemente enthält und unter den Regeln abgeschlossen ist.
\end{definitionbox}

"`Kleinste"' ist eine Rechenvorschrift: Regeln anwenden, bis nichts Neues
mehr entsteht. Formal ist die definierte Menge der kleinste Fixpunkt des
zugehörigen Operators, was im Kapitel über Verbände wieder auftaucht.

\begin{examplebox}
Terme über zwei Konstanten und einer zweistelligen Operation:

\begin{center}
\begin{tabular}{lrrr}
\toprule
Schachtelungstiefe & 0 & 1 & 2\\
\midrule
Anzahl Terme & 2 & 6 & 38\\
\bottomrule
\end{tabular}
\end{center}

Die Rekursion ist $t(n+1) = c + t(n)^2$ mit $c$ der Anzahl der Konstanten.
Der naheliegende Fehler ist $t(n) + t(n)^2 = 42$; dabei werden die vier
bereits vorhandenen Terme der Tiefe 1 doppelt gezählt.
\end{examplebox}

\subsection{Induktives Beweisen}

\begin{definitionbox}
\textbf{Vollständige Induktion.} Gilt $P(0)$ und folgt aus $P(n)$ stets
$P(n+1)$, so gilt $P(n)$ für alle $n \in \mathbb{N}$.

\textbf{Strukturelle Induktion.} Gilt eine Eigenschaft für die Basisfälle und
überträgt sie sich bei jedem Konstruktionsschritt, so gilt sie für alle
Elemente der induktiv definierten Menge.
\end{definitionbox}

\begin{notebox}
Das Aufzählen aller Strukturen bis zu einer Tiefe ersetzt den Beweis nicht.
Es leistet aber das, was der Beweis nicht leistet: bei einer falschen
Behauptung liefert es sofort ein Gegenbeispiel. Die meisten gescheiterten
Beweisversuche beginnen als Behauptungen dieser Art.
\end{notebox}

\subsection{Fundierte Ordnungen}

\begin{definitionbox}
Eine Relation heißt \emph{fundiert} (noethersch), wenn es keine unendliche
absteigende Kette gibt. Über einer endlichen Menge ist das gleichbedeutend
mit Zyklenfreiheit.
\end{definitionbox}

Drei Terminierungsargumente, die dasselbe sind:
\begin{enumerate}
  \item ein \emph{Maß} nach $\mathbb{N}$, das bei jedem Schritt echt fällt;
  \item eine \emph{lexikographische} Ordnung, etwa für die Ackermannfunktion,
        bei der kein einzelnes Argument fällt, das Paar aber schon;
  \item \emph{noethersche Induktion}: eine Eigenschaft, die initial gilt und
        bei jedem Schritt erhalten bleibt, gilt überall im erreichbaren
        Zustandsraum.
\end{enumerate}

\subsection{Rekursionssatz}

\begin{definitionbox}
Ist $M$ induktiv durch Basiselemente und Regeln definiert und sind Werte für
die Basiselemente sowie eine Vorschrift für jede Regel gegeben, so existiert
genau eine Funktion auf $M$, die diesen Gleichungen genügt.
\end{definitionbox}

Das ist die Rechtfertigung dafür, Funktionen durch Rekursionsgleichungen zu
definieren: die Existenz einer solchen Funktion ist ein Satz und keine
Selbstverständlichkeit. Sie scheitert, wenn die Definition nicht fundiert
ist, etwa bei $f(n) = f(n+1)$.

\section{Ordnungsstrukturen}

\subsection{Halbordnungen}

\begin{definitionbox}
$(M, \sqsubseteq)$ heißt \emph{Halbordnung}, wenn $\sqsubseteq$ reflexiv,
antisymmetrisch und transitiv ist. Sie heißt \emph{total}, wenn je zwei
Elemente vergleichbar sind.
\end{definitionbox}

Das Hasse-Diagramm zeichnet nur die Überdeckungsrelation, also die Paare, bei
denen nichts echt dazwischen liegt. Aus ihr lässt sich die Ordnung als
transitive Hülle zurückgewinnen, weshalb nichts verloren geht.

\begin{examplebox}
Die Teilmengen einer dreielementigen Menge unter $\subseteq$:

\begin{center}
\begin{tabular}{lr}
\toprule
Elemente & 8\\
längste Kette & 4\\
breiteste Antikette & 3\\
lineare Erweiterungen & 48\\
\bottomrule
\end{tabular}
\end{center}

Die 48 linearen Erweiterungen sind genau die zulässigen topologischen
Sortierungen. Eine totale Ordnung hat genau eine; deshalb ist eine
topologische Sortierung eine Auswahl und keine Berechnung.
\end{examplebox}

\subsection{Verbände}

\begin{definitionbox}
Eine Halbordnung, in der je zwei Elemente ein Supremum $a \sqcup b$ und ein
Infimum $a \sqcap b$ besitzen, heißt \emph{Verband}.
\end{definitionbox}

Äquivalent ist die algebraische Sicht: eine Menge mit zwei Operationen, die
idempotent, kommutativ und assoziativ sind und das \emph{Absorptionsgesetz}
\[
  a \sqcap (a \sqcup b) = a, \qquad a \sqcup (a \sqcap b) = a
\]
erfüllen. Die Absorption ist das Gesetz, das die beiden Operationen
verbindet; ohne sie wären sie unabhängig voneinander, und aus ihr gewinnt man
die Ordnung zurück als $a \sqsubseteq b \iff a \sqcup b = b$.

\subsubsection{Distributiv und modular}

\begin{definitionbox}
Ein Verband heißt \emph{distributiv}, wenn
$a \sqcap (b \sqcup c) = (a \sqcap b) \sqcup (a \sqcap c)$ gilt, und
\emph{modular}, wenn diese Gleichung wenigstens für $a \sqsubseteq c$ gilt.
\end{definitionbox}

\begin{center}
\begin{tabular}{lcc}
\toprule
& modular & distributiv\\
\midrule
$M_3$ (Diamant) & ja & nein\\
$N_5$ (Pentagon) & nein & nein\\
\bottomrule
\end{tabular}
\end{center}

Diese beiden fünfelementigen Verbände beantworten beide Fragen vollständig:
ein Verband ist modular genau dann, wenn er $N_5$ nicht enthält, und
distributiv genau dann, wenn er weder $N_5$ noch $M_3$ enthält.

\subsection{Vollständige Verbände}

\begin{definitionbox}
Ein Verband heißt \emph{vollständig}, wenn jede Teilmenge ein Supremum und
ein Infimum besitzt.
\end{definitionbox}

Jeder endliche Verband ist vollständig. Der Fall der leeren Teilmenge sieht
wie eine Ausnahme aus und ist keine: jedes Element ist trivialerweise obere
Schranke der leeren Menge, also ist $\bigsqcup \emptyset$ das kleinste
Element und $\bigsqcap \emptyset$ das größte.

\subsection{Boolesche Verbände}

\begin{definitionbox}
Ein \emph{Boolescher Verband} ist ein distributiver Verband mit $0$ und $1$,
in dem jedes Element ein Komplement besitzt.
\end{definitionbox}

\begin{examplebox}
Der Teilerverband von $30 = 2 \cdot 3 \cdot 5$ ist Boolesch, der von
$12 = 2^2 \cdot 3$ nicht. Bei $12$ bilden $1, 2, 4$ eine Kette, und $2$ hat
kein Komplement.
\end{examplebox}

\textbf{Darstellungssatz (endlicher Fall).} Die Abbildung, die jedem Element
die Menge der unter ihm liegenden Atome zuordnet, ist ein Isomorphismus auf
die Potenzmenge der Atome. Folglich hat jeder endliche Boolesche Verband
$2^k$ Elemente, und es gibt keinen mit sechs oder zwölf Elementen.

\subsection{Fixpunkte}

\begin{definitionbox}
\textbf{Satz von Knaster und Tarski.} Ist $L$ ein vollständiger Verband und
$f: L \to L$ monoton, so besitzt $f$ einen kleinsten und einen größten
Fixpunkt, nämlich
\[
  \mathrm{lfp}(f) = \bigsqcap \{ x \mid f(x) \sqsubseteq x \}, \qquad
  \mathrm{gfp}(f) = \bigsqcup \{ x \mid x \sqsubseteq f(x) \}.
\]
\end{definitionbox}

Auf einem endlichen Verband erreicht auch die Kleene-Iteration
$\bot, f(\bot), f^2(\bot), \dots$ den kleinsten Fixpunkt, und zwar nach
höchstens so vielen Schritten, wie der Verband hoch ist. Im Begleitcode:
4 Schritte bei einer Höhe von 5.

\begin{notebox}
Auf unendlichen Verbänden fallen die beiden Konstruktionen auseinander: die
Iteration braucht Stetigkeit, nicht nur Monotonie. Genau deshalb lässt sich
eine Datenflussanalyse über einem endlichen Verband durch Iterieren
berechnen und über einem unendlichen nicht.
\end{notebox}

\subsection{Strukturverträgliche Abbildungen}

\begin{definitionbox}
$f$ heißt \emph{Ordnungshomomorphismus} (monoton), wenn aus
$a \sqsubseteq b$ stets $f(a) \sqsubseteq f(b)$ folgt, und
\emph{Verbandshomomorphismus}, wenn $f(a \sqcup b) = f(a) \sqcup f(b)$ und
$f(a \sqcap b) = f(a) \sqcap f(b)$ gelten.
\end{definitionbox}

Jeder Verbandshomomorphismus ist monoton, denn die Ordnung ist aus jeder der
beiden Operationen definierbar. Die Umkehrung gilt nicht: eine monotone
Abbildung darf zwei unvergleichbare Elemente auf dasselbe Bild werfen, und
dann hat ihr Supremum keinen passenden Platz mehr.

\subsection{Konstruktionsprinzipien}

Aus vorhandenen Verbänden entstehen neue:
\begin{itemize}
  \item \textbf{Unterverband}: eine unter $\sqcup$ und $\sqcap$
        abgeschlossene Teilmenge. Achtung: eine Teilmenge, die als Ordnung
        ein Verband ist, muss kein Unterverband sein, weil Supremum und
        Infimum in der Teilmenge andere sein können als im Oberverband.
  \item \textbf{Produkt}: $L_1 \times L_2$ mit komponentenweiser Ordnung.
        Es ist genau dann distributiv, wenn beide Faktoren es sind.
  \item \textbf{Dual}: dieselbe Menge mit umgekehrter Ordnung. Jede Aussage
        über Verbände hat eine duale Aussage, in der $\sqcup$ und $\sqcap$
        vertauscht sind, und diese gilt automatisch mit.
\end{itemize}

\begin{notebox}
Das Dualitätsprinzip halbiert die Arbeit: wer die Absorption in einer Form
gezeigt hat, hat sie in beiden gezeigt. Es gilt, weil die Axiome selbst
dual-symmetrisch sind.
\end{notebox}

\section{Algebraische Strukturen}

\subsection{Eine Verknüpfung}

\begin{center}
\begin{tabular}{lll}
\toprule
Struktur & Bedingungen\\
\midrule
Magma & abgeschlossen\\
Halbgruppe & $+$ assoziativ\\
Monoid & $+$ neutrales Element\\
Gruppe & $+$ Inverse\\
\bottomrule
\end{tabular}
\end{center}

\begin{examplebox}
Wie stark die Assoziativität einschränkt, gezählt im Begleitcode:

\begin{center}
\begin{tabular}{lrr}
\toprule
Trägermenge & Verknüpfungen & davon assoziativ\\
\midrule
2 Elemente & 16 & 8\\
3 Elemente & 19\,683 & 113\\
\bottomrule
\end{tabular}
\end{center}
\end{examplebox}

Das neutrale Element ist eindeutig: gäbe es zwei, so ließe jedes das andere
unverändert, also wären sie gleich.

\subsection{Gruppen}

\begin{definitionbox}
$(G, \cdot)$ ist eine \emph{Gruppe}, wenn die Verknüpfung abgeschlossen und
assoziativ ist, ein neutrales Element existiert und jedes Element ein
Inverses besitzt. Gilt zusätzlich $ab = ba$, heißt die Gruppe
\emph{abelsch}.
\end{definitionbox}

Die Invertierbarkeit trägt die ganze Theorie: sie liefert die Kürzungsregel,
macht jede Zeile der Verknüpfungstafel zu einer Permutation und erzwingt,
dass die Ordnung jedes Elements die Gruppenordnung teilt.

\begin{examplebox}
Die fünf Gruppen der Ordnung 8, im Begleitcode durch Suche über alle
Bijektionen klassifiziert:

\begin{center}
\begin{tabular}{ll}
\toprule
Gruppe & abelsch\\
\midrule
$\mathbb{Z}_8$ & ja\\
$\mathbb{Z}_4 \times \mathbb{Z}_2$ & ja\\
$\mathbb{Z}_2 \times \mathbb{Z}_2 \times \mathbb{Z}_2$ & ja\\
$D_4$ & nein\\
$Q_8$ & nein\\
\bottomrule
\end{tabular}
\end{center}
\end{examplebox}

\subsection{Untergruppen und Nebenklassen}

\begin{definitionbox}
\textbf{Satz von Lagrange.} Ist $U \leq G$ endlich, so teilt $|U|$ die
Gruppenordnung $|G|$, und $[G:U] \cdot |U| = |G|$.
\end{definitionbox}

Der Beweis ist ein Abzählargument: die Linksnebenklassen $aU$ zerlegen $G$
und haben alle die Größe $|U|$. Beide Hälften sind im Begleitcode für alle
Untergruppen mehrerer Gruppen nachgerechnet.

\begin{center}
\begin{tabular}{lrrr}
\toprule
Gruppe & Ordnung & Untergruppen & davon normal\\
\midrule
$\mathbb{Z}_{12}$ & 12 & 6 & 6\\
$S_3$ & 6 & 6 & 3\\
$D_4$ & 8 & 10 & 6\\
$Q_8$ & 8 & 6 & 6\\
$\mathbb{Z}_2 \times \mathbb{Z}_2$ & 4 & 5 & 5\\
\bottomrule
\end{tabular}
\end{center}

\begin{notebox}
$Q_8$ ist nicht abelsch, und trotzdem sind alle sechs Untergruppen normal.
Normalität ist also \emph{keine} abgeschwächte Kommutativität. In $S_3$
dagegen sind genau die drei Untergruppen der Ordnung 2 nicht normal, und das
sind genau die, bei denen Links- und Rechtsnebenklassen auseinanderfallen.
\end{notebox}

\subsection{Faktorgruppen}

\begin{definitionbox}
$N \trianglelefteq G$ heißt \emph{Normalteiler}, wenn $gNg^{-1} = N$ für alle
$g \in G$ gilt. Dann ist $G/N$ mit $(aN)(bN) = (ab)N$ eine Gruppe.
\end{definitionbox}

Der einzige Punkt, der zu prüfen ist, ist die \emph{Wohldefiniertheit}: das
Produkt darf nicht von der Wahl der Repräsentanten abhängen. Genau das ist
die Normalität, und bei einer nicht normalen Untergruppe von $S_3$ liefern
zwei Repräsentantenwahlen Ergebnisse in verschiedenen Nebenklassen.

\subsection{Homomorphismen}

\begin{definitionbox}
\textbf{Homomorphiesatz.} Für einen Gruppenhomomorphismus
$\varphi: G_1 \to G_2$ gilt: $\ker \varphi \trianglelefteq G_1$,
$\operatorname{im} \varphi \leq G_2$, und
\[
  G_1 / \ker \varphi \;\cong\; \operatorname{im} \varphi .
\]
\end{definitionbox}

Der Kern misst genau, was die Abbildung verliert: zwei Elemente haben dasselbe
Bild genau dann, wenn sie sich um ein Kernelement unterscheiden. Injektivität
ist der Fall $\ker \varphi = \{e\}$.

\subsection{Ringe, Ideale, Körper}

\begin{definitionbox}
Ein \emph{Ring} ist eine abelsche Gruppe bezüglich $+$ und eine Halbgruppe
bezüglich $\cdot$, verbunden durch die Distributivgesetze. Ein \emph{Ideal}
$I \trianglelefteq R$ ist eine additive Untergruppe mit $rI \subseteq I$ und
$Ir \subseteq I$ für alle $r \in R$.
\end{definitionbox}

Das Ideal spielt für Ringe die Rolle, die der Normalteiler für Gruppen
spielt: es ist genau die Bedingung, unter der $R/I$ wieder ein Ring ist.

\begin{examplebox}
$\mathbb{Z}/12\mathbb{Z}$ hat sechs Ideale, und 12 hat sechs Teiler. Jedes
Ideal ist ein Hauptideal, im Begleitcode für alle sechs nachgeprüft.
Einheiten sind $1, 5, 7, 11$; Nullteiler sind $2, 3, 4, 6, 8, 9, 10$.
\end{examplebox}

\begin{definitionbox}
Ein \emph{Integritätsbereich} ist ein kommutativer Ring mit Eins ohne
Nullteiler. Ein \emph{Körper} ist ein Integritätsbereich, in dem jedes
Element $\neq 0$ invertierbar ist.
\end{definitionbox}

$\mathbb{Z}/n\mathbb{Z}$ ist genau dann ein Körper, wenn $n$ prim ist; im
Begleitcode für alle $n < 30$ geprüft. Der Grund ist immer derselbe: ein
zusammengesetzter Modul hat Nullteiler, und ein Nullteiler kann nicht
invertierbar sein.

\begin{notebox}
Das sagt nichts darüber, welche \emph{Größen} ein Körper haben kann. Es gibt
einen Körper mit vier Elementen, und er ist nicht $\mathbb{Z}/4\mathbb{Z}$,
sondern $\mathbb{F}_2[x]/(x^2+x+1)$. Seine Charakteristik ist 2, nicht 4.
Deshalb werden endliche Körper nach Primzahl\emph{potenzen} klassifiziert.
\end{notebox}

\subsection{Untergruppenkriterium}

\begin{definitionbox}
$U \subseteq G$ ist genau dann Untergruppe, wenn $U \neq \emptyset$ und
für alle $a, b \in U$ auch $ab^{-1} \in U$ gilt.
\end{definitionbox}

Für endliches $G$ genügt sogar die Abgeschlossenheit: ist $U$ endlich,
nichtleer und unter der Verknüpfung abgeschlossen, so enthält $U$ mit $a$
auch alle Potenzen $a^k$, und da diese sich wiederholen müssen, liegt das
neutrale Element und damit das Inverse in $U$. Für unendliche Gruppen ist das
falsch: $\mathbb{N}$ ist unter Addition abgeschlossen und keine Untergruppe
von $\mathbb{Z}$.

\subsection{Schnitte und Produkte}

Der Schnitt beliebig vieler Untergruppen ist wieder eine Untergruppe; die
Vereinigung im Allgemeinen nicht. Deshalb ist "`die von $S$ erzeugte
Untergruppe"' wohldefiniert: sie ist der Schnitt aller Untergruppen, die $S$
enthalten.

Dasselbe Muster gilt für Ideale, für Unterverbände und für Untervektorräume,
und es ist derselbe Grund: die definierende Eigenschaft ist unter Durchschnitt
stabil. Daraus entsteht jedes Mal ein Hüllenoperator, und damit wieder ein
kleinster Fixpunkt im Sinne des Satzes von Tarski.

\section{Lineare Algebra}

\subsection{Vektorräume}

\begin{definitionbox}
Eine Familie $(v_1, \dots, v_n)$ heißt \emph{linear unabhängig}, wenn
$\sum \lambda_i v_i = 0$ nur mit $\lambda_1 = \dots = \lambda_n = 0$ möglich
ist. Eine \emph{Basis} ist eine unabhängige Familie, die den Raum aufspannt;
die \emph{Dimension} ist ihre Länge.
\end{definitionbox}

\textbf{Steinitzscher Austauschsatz.} Ein Vektor aus dem Erzeugnis einer
Familie kann ein Mitglied ersetzen, ohne das Erzeugnis zu ändern. Wiederholtes
Austauschen zeigt, dass je zwei Basen gleich lang sind, und erst dadurch ist
die Dimension eine Eigenschaft des Raums und nicht der gewählten Basis.

Alle vier Begriffe reduzieren sich auf den Rang einer Matrix: eine Familie ist
unabhängig, wenn ihr Rang ihrer Länge entspricht, ein Vektor liegt im
Erzeugnis, wenn Hinzufügen den Rang nicht erhöht, und die Dimension ist der
Rang.

\subsection{Das Gaußsche Eliminationsverfahren}

\begin{definitionbox}
\textbf{Hauptsatz über lineare Gleichungssysteme.} Für $Ax = b$ gilt:
\begin{itemize}
  \item $\operatorname{rg}(A) < \operatorname{rg}(A|b)$: keine Lösung;
  \item $\operatorname{rg}(A) = \operatorname{rg}(A|b) < n$: unendlich viele;
  \item $\operatorname{rg}(A) = \operatorname{rg}(A|b) = n$: genau eine.
\end{itemize}
\end{definitionbox}

Die Lösungsmenge ist eine spezielle Lösung plus der Kern, also ein Punkt und
die Richtungen, in denen man sich bewegen darf. Die Stufenform hängt von der
Reihenfolge der Zeilen ab, die \emph{reduzierte} Stufenform nicht: sie ist
eindeutig, weshalb sie entscheidet, ob zwei Systeme dieselbe Lösungsmenge
haben.

\begin{notebox}
Im Begleitcode ist jeder Eintrag ein exakter Bruch. Ein Pivot ist genau dann
null, wenn er wirklich null ist. In Gleitkomma muss dieselbe Rechnung raten,
ob eine kleine Zahl ein Rundungsfehler oder ein Wert ist, und dieses Raten
ist der ganze Unterschied zwischen numerischer linearer Algebra und der
Version, über die die Vorlesung Sätze beweist.
\end{notebox}

\subsection{Invertieren von Matrizen}

Gauß-Jordan neben der Einheitsmatrix: dieselben Zeilenoperationen wirken auf
beide Hälften, also steht rechts am Ende die Inverse. Erreicht die linke
Hälfte die Einheitsmatrix nicht, ist die Matrix singulär.

Jede Zeilenoperation ist selbst eine Matrix. Wendet die Elimination
$E_1, \dots, E_k$ an, so ist
\[
  E_k \cdots E_1 A = I \quad\Longrightarrow\quad
  A = E_1^{-1} E_2^{-1} \cdots E_k^{-1}.
\]
Die Elimination \emph{ist} die Faktorisierung: die invertierbaren Matrizen
werden von den Elementarmatrizen erzeugt.

\subsection{Lineare Abbildungen}

\begin{definitionbox}
Eine lineare Abbildung ist durch die Bilder einer Basis festgelegt. Die
Matrix einer Abbildung bezüglich zweier Basen hat als Spalten die
\emph{Koordinaten} der Bilder der Quellbasis in der Zielbasis.
\end{definitionbox}

Ein Basiswechsel ändert die Matrix, nicht die Abbildung. Was invariant
bleibt, beschreibt die Abbildung selbst: Rang, Determinante, Eigenwerte.

\begin{definitionbox}
\textbf{Dimensionssatz.} $\dim \ker f + \dim \operatorname{im} f =
\dim V$ für $f: V \to W$.
\end{definitionbox}

Daraus folgt: eine Abbildung zwischen Räumen gleicher Dimension ist genau
dann injektiv, wenn sie surjektiv ist, und eine Abbildung in einen Raum
kleinerer Dimension ist nie injektiv. Der Homomorphiesatz der Gruppentheorie
ist derselbe Satz: $V/\ker f \cong \operatorname{im} f$.

\subsection{Determinanten}

Drei Definitionen, eine Zahl, sehr verschiedene Kosten (Multiplikationen):

\begin{center}
\begin{tabular}{lrrr}
\toprule
Größe & Leibniz & Laplace & Gauß\\
\midrule
4 & 72 & 40 & 20\\
6 & 3\,600 & 1\,236 & 70\\
8 & 282\,240 & 69\,280 & 168\\
\bottomrule
\end{tabular}
\end{center}

Alle drei sind im Begleitcode implementiert und stimmen überein. Das ist der
Sinn davon, auch die teuren zu implementieren: die billige Methode wird
geglaubt, weil sie mit der Definition übereinstimmt, nicht weil die
Definition unbrauchbar wäre.

Wesentliche Eigenschaften: $\det(AB) = \det(A)\det(B)$; $\det A \neq 0$ genau
dann, wenn $A$ invertierbar ist; ein Zeilentausch wechselt das Vorzeichen.

\subsection{Eigenwerte}

\begin{definitionbox}
$\lambda$ heißt \emph{Eigenwert} von $A$, wenn $Av = \lambda v$ für ein
$v \neq 0$ gilt. Die Eigenwerte sind die Nullstellen des charakteristischen
Polynoms $\chi_A(x) = \det(A - xI)$.
\end{definitionbox}

Für $2 \times 2$-Matrizen ist $\chi_A(x) = x^2 - \operatorname{tr}(A) x +
\det(A)$.

Zwei verschiedene Arten zu scheitern:
\begin{itemize}
  \item eine Drehung hat \emph{keinen} reellen Eigenwert, weil keine Richtung
        erhalten bleibt;
  \item eine Scherung $\begin{pmatrix}1&1\\0&1\end{pmatrix}$ hat einen
        Eigenwert und zu wenige Eigenvektoren, ist also nicht
        diagonalisierbar, obwohl ihr Eigenwert reell ist.
\end{itemize}

Diagonalisierbarkeit ist eine Aussage über Eigenvektoren, nicht über die
Existenz von Eigenwerten.

\subsection{Anwendung: Eigenfaces}

Ein Gesicht ist ein Punkt in einem Raum mit einer Dimension pro Pixel. Die
Hauptkomponenten sind die Richtungen der größten Varianz, also die
Eigenvektoren der Kovarianzmatrix.

\begin{examplebox}
Acht synthetische Bilder, aus drei Grundmustern erzeugt:

\begin{center}
\begin{tabular}{lr}
\toprule
Komponenten & Rekonstruktionsfehler\\
\midrule
1 & 0{,}0243\\
2 & 0{,}0073\\
3 & 0{,}00003\\
8 & 0\\
\bottomrule
\end{tabular}
\end{center}

Der Fehler bricht bei drei Komponenten ein, und die erklärte Varianz sagt
dasselbe: 0{,}76, 0{,}17, 0{,}07, dann 0{,}0001. Die Methode findet die drei
Muster, aus denen die Daten gebaut wurden.
\end{examplebox}

Die gesamte Pipeline besteht aus Bausteinen der vorigen Abschnitte: eine
orthonormale Basis, eine Projektion als lineare Abbildung, Koordinaten
bezüglich dieser Basis und eine Nächste-Nachbarn-Suche in diesen Koordinaten.

\subsection{Algebra der linearen Abbildungen}

Die linearen Abbildungen $V \to V$ bilden mit punktweiser Addition und
Komposition einen Ring, dessen Einheiten genau die Isomorphismen sind. Unter
Wahl einer Basis ist dieser Ring isomorph zum Matrizenring, und genau das ist
der Grund, warum Matrizenrechnung Abbildungsrechnung ist.

\begin{definitionbox}
\textbf{Basiswechsel.} Sind $B$ und $B'$ Basen von $V$ und ist $S$ die
Matrix des Basiswechsels, so gilt für die Matrizen $A$ und $A'$ derselben
Abbildung
\[
  A' = S^{-1} A S .
\]
\end{definitionbox}

Zwei Matrizen, die durch eine solche Konjugation auseinander hervorgehen,
heißen \emph{ähnlich}. Ähnliche Matrizen haben denselben Rang, dieselbe
Determinante, dieselbe Spur und dasselbe charakteristische Polynom, weil all
diese Größen Eigenschaften der Abbildung sind und nicht der Basis.

\subsection{Ein durchgerechnetes Beispiel}

\[
  A = \begin{pmatrix} 2 & 1 \\ 1 & 1 \end{pmatrix}, \quad
  b = \begin{pmatrix} 5 \\ 5 \end{pmatrix}
\]

Elimination: die zweite Zeile minus $\tfrac12$ der ersten ergibt
$\bigl(0, \tfrac12 \mid \tfrac52\bigr)$, also $x_2 = 5$ (nach
Rückeinsetzen: $x_2 = 1$, $x_1 = 2$ mit der zweiten Normierung). Der Rang ist
2, also gleich der Anzahl der Unbekannten: genau eine Lösung, nämlich
$(2, 1)$.

Determinante: $\det A = 2 \cdot 1 - 1 \cdot 1 = 1 \neq 0$, also ist $A$
invertierbar, und
\[
  A^{-1} = \begin{pmatrix} 1 & -1 \\ -1 & 2 \end{pmatrix}.
\]

Charakteristisches Polynom: $x^2 - 3x + 1$, mit den Nullstellen
$\tfrac{3 \pm \sqrt5}{2}$. Beide sind irrational, also findet die exakte
Suche über rationale Nullstellen im Begleitcode hier nichts, was der korrekte
Befund ist und nicht ein Fehlschlag: die Eigenwerte existieren, sie sind nur
nicht rational.

\section{Zusammenfassung}

Vier Ergebnisse, die der ersten Vermutung widersprechen:

\begin{enumerate}
  \item \textbf{Normalität ist keine abgeschwächte Kommutativität.} $Q_8$ ist
        nicht abelsch, und alle sechs Untergruppen sind normal.
  \item \textbf{Es gibt einen Körper mit vier Elementen, und er ist nicht
        $\mathbb{Z}/4\mathbb{Z}$.} Die Aussage "`$\mathbb{Z}/n\mathbb{Z}$ ist
        genau für Primzahlen ein Körper"' sagt nichts über mögliche Größen.
  \item \textbf{Hüllen vertauschen nicht.} Erst symmetrisch, dann transitiv
        ergibt eine Äquivalenzrelation; die andere Reihenfolge ergibt eine
        Relation, die nicht einmal transitiv ist.
  \item \textbf{Die Termanzahl ist 38, nicht 42.} Die Rekursion lautet
        $c + t(n)^2$, denn bereits gebaute Terme werden nicht neu gebaut.
\end{enumerate}

Der Begleitcode liegt in \texttt{college-assignment/math1/} und besteht aus
32 Modulen in sechs Blöcken. Alle Zahlen dieses Skripts stammen aus diesen
Modulen; die lineare Algebra rechnet exakt in Brüchen und ist gegen numpy auf
800 zufälligen Eingaben geprüft.

\end{document}
